\chapter{presentation monologue}

\section{Introduction}
hello everyone, today I will present my master thesis on comparing 
two methods for adjusting for time-invariant unmearued confounding in the context of 
longitudinal causal inference with time varying treatments. The two methods I will compare are fixed effects models 
introduced by \cite{blackwellEffectPoliticalAdvertising2024} and sequenttial guassian process latent variable models introduced by \cite{hattSequentialDeconfoundingCausal2024}.

\section{simple scenario}
To help you conceptualize the problem, imagine we have a simple scenario where we have a blood pressure pill, and
 10000000 individuals, 12 months, and at the end of this 12 months, we are going to measure the blooad pressure of these individuals, 
and we want to see if there is a difference in the blood pressure between those who took the pill and those who did not take the pill in the 
last month with the same treament pattern (meaning they took the pill in the same months) in the past 11 months. 
The best way to answer this question is to run a randomized 
controlled trial, in which at each month, we randommly assign individuals to either take the pill or not take the pill, 
and then we measure the outcome at the end of the 12th month. Then for each treatment pattern in the past 11 months, 
we can compare the outcome between those who took the pill in the last month and those who did not take the pill in the last month. This 
difference is the causal effect of taking the pill in the last month, given the treatment pattern in the past 11 months. This is just an example of 
different questions we could ask in such a study.  

But usually in practice, we can't run such a study, becuase of ethical, financial, or logistical reasons,
and we have to rely on observational data, meaning we observe the treatment and outcome without random assignment,
and in such a case, we have to deal with confounding. Counfounding is a situation where there is a third variable 
that is associated with both the treatment and the outcome, and this can bias our estimates of the causal effect of the treatment on the outcome if 
we do not adjust for it. If we have measured all the confounders, we are in a situtaion called "no unmeasured confounding", 
and in such a case, we can adjust for them using statistical methods such as inverse probability weighting.
If we have measured some of the confounders, but not all of them, we are in a situation called "unmeasured confounding". 
in a longitudinal study, we can have two types of unmeasured confounding: time-varying unmeasured confounding and time-invariant unmeasured confounding.
In this thesis, we focus on time-invariant unmeasured confounding, which is a situation where there is a 
confounder that does not change over time, but it is not measured in our data. In our blood pressure example, this could be something like 
genetic predisposition to high blood pressure, which does not change over time, but it is not measured in our data.

\section{methods}
There are different methods to tackle time-invariant unmeasured confounding, and here, I focus on two methods that are 
close to each other conceptually, but they are different in terms of their assumptions and implementation. These two methods are: 
fixed effects models and sequential gaussian process latent variable models. The idea behind both of these methods is 
to use the treatment history of individuals to retrieve the foot print of the unmeasured confounder, 
and then use this footprint to adjust for the confounding bias. By doing this, we regain the sequential ignorability assumption, 
which says the potential outcomes of at each time point for a given assigned treatments until that time point 
are independent of the treatment assignment at that time point, conditioned on the treatment history, observed confounders, 
and the unmeaured confounder substitute. 


in the past are independent of the treatment assignment at that time point, given the treatment history in the past.



a substitute value for the unmeasured confounder, which gives 
 meaning it does not change over time, and therefore it leaves a consistent footprint in the treatment history of individuals.








